\documentclass[12pt,a4paper]{article}

% ============================
% PACKAGES
% ============================
\usepackage[utf8]{inputenc}
\usepackage[T1]{fontenc}
\usepackage{geometry}
\usepackage{graphicx}
\usepackage{xcolor}
\usepackage{titlesec}
\usepackage{tocloft}
\usepackage{hyperref}
\usepackage{enumitem}
\usepackage{booktabs}
\usepackage{array}
\usepackage{longtable}
\usepackage{fancyhdr}
\usepackage{tikz}
\usepackage{tcolorbox}
\usepackage{fontawesome5}
\usepackage{setspace}
\usepackage{parskip}

% ============================
% PAGE GEOMETRY
% ============================
\geometry{
    left=2.5cm,
    right=2.5cm,
    top=2.5cm,
    bottom=2.5cm
}

% ============================
% COLORS
% ============================
\definecolor{primaryblue}{RGB}{0, 82, 147}
\definecolor{secondaryblue}{RGB}{52, 152, 219}
\definecolor{lightblue}{RGB}{214, 234, 248}
\definecolor{darkgray}{RGB}{52, 73, 94}
\definecolor{lightgray}{RGB}{236, 240, 241}
\definecolor{accentgreen}{RGB}{39, 174, 96}
\definecolor{accentorange}{RGB}{230, 126, 34}
\definecolor{accentred}{RGB}{192, 57, 43}

% ============================
% HYPERREF SETUP
% ============================
\hypersetup{
    colorlinks=true,
    linkcolor=primaryblue,
    urlcolor=secondaryblue,
    citecolor=primaryblue,
    pdfauthor={Business Studies},
    pdftitle={Comprehensive Business Studies Booklet},
    pdfsubject={Business Education}
}

% ============================
% HEADER AND FOOTER
% ============================
\pagestyle{fancy}
\fancyhf{}
\fancyhead[L]{\textcolor{primaryblue}{\textbf{Business Studies Booklet}}}
\fancyhead[R]{\textcolor{darkgray}{\thepage}}
\fancyfoot[C]{\textcolor{darkgray}{\small A Comprehensive Guide to Business Fundamentals}}
\renewcommand{\headrulewidth}{2pt}
\renewcommand{\headrule}{\hbox to\headwidth{\color{primaryblue}\leaders\hrule height \headrulewidth\hfill}}

% ============================
% SECTION FORMATTING
% ============================
\titleformat{\section}
    {\color{primaryblue}\LARGE\bfseries}
    {\thesection}{1em}{}[\titlerule]
    
\titleformat{\subsection}
    {\color{secondaryblue}\Large\bfseries}
    {\thesubsection}{1em}{}
    
\titleformat{\subsubsection}
    {\color{darkgray}\large\bfseries}
    {\thesubsubsection}{1em}{}

% ============================
% TOC FORMATTING
% ============================
\renewcommand{\cftsecfont}{\color{primaryblue}\bfseries}
\renewcommand{\cftsubsecfont}{\color{secondaryblue}}
\renewcommand{\cftsecpagefont}{\color{primaryblue}\bfseries}

% ============================
% TCOLORBOX STYLES
% ============================
\tcbuselibrary{skins, breakable}

\newtcolorbox{definitionbox}[1][]{
    enhanced,
    colback=lightblue,
    colframe=primaryblue,
    coltitle=white,
    fonttitle=\bfseries,
    title=\faIcon{book} Definition,
    #1
}

\newtcolorbox{keyconceptbox}[1][]{
    enhanced,
    colback=lightgray,
    colframe=darkgray,
    coltitle=white,
    fonttitle=\bfseries,
    title=\faIcon{lightbulb} Key Concept,
    #1
}

\newtcolorbox{examplebox}[1][]{
    enhanced,
    colback=green!5,
    colframe=accentgreen,
    coltitle=white,
    fonttitle=\bfseries,
    title=\faIcon{check-circle} Example,
    #1
}

\newtcolorbox{importantbox}[1][]{
    enhanced,
    colback=orange!5,
    colframe=accentorange,
    coltitle=white,
    fonttitle=\bfseries,
    title=\faIcon{exclamation-triangle} Important,
    #1
}

\newtcolorbox{advantagebox}[1][]{
    enhanced,
    colback=green!10,
    colframe=accentgreen,
    coltitle=white,
    fonttitle=\bfseries,
    title=\faIcon{thumbs-up} Advantages,
    #1
}

\newtcolorbox{disadvantagebox}[1][]{
    enhanced,
    colback=red!10,
    colframe=accentred,
    coltitle=white,
    fonttitle=\bfseries,
    title=\faIcon{thumbs-down} Disadvantages,
    #1
}

% ============================
% DOCUMENT BEGINS
% ============================
\begin{document}

% ============================
% TITLE PAGE
% ============================
\begin{titlepage}
    \centering
    \vspace*{2cm}
    
    {\Huge\bfseries\textcolor{primaryblue}{Business Studies}}\\[0.5cm]
    {\Huge\bfseries\textcolor{secondaryblue}{Comprehensive Booklet}}\\[1cm]
    
    \begin{tikzpicture}
        \draw[line width=3pt, primaryblue] (0,0) -- (12,0);
    \end{tikzpicture}
    
    \vspace{1.5cm}
    
    {\Large\textcolor{darkgray}{A Complete Guide to Understanding}}\\[0.3cm]
    {\Large\textcolor{darkgray}{Business Fundamentals \& Concepts}}
    
    \vspace{2cm}
    
    \begin{tikzpicture}
        \fill[primaryblue] (0,0) circle (2.5cm);
        \node[white] at (0,0) {\Huge\faIcon{briefcase}};
    \end{tikzpicture}
    
    \vfill
    
    {\large\textcolor{darkgray}{\today}}
    
\end{titlepage}

% ============================
% TABLE OF CONTENTS
% ============================
\newpage
\tableofcontents
\newpage

% ============================
% CHAPTER 1: INTRODUCTION TO BUSINESS STUDIES
% ============================
\section{Introduction to Business Studies}

\begin{definitionbox}
\textbf{Business Studies} is an academic subject taught in schools and at university level in many countries. Its study combines elements of accountancy, finance, marketing, organizational studies, human resource management, and operations. Business studies is a broad subject, where the range of topics is designed to give the student a general overview of the various elements of running a business.
\end{definitionbox}

\subsection{What Is a Business?}

\begin{keyconceptbox}
The term \textbf{business} refers to an organization or enterprising entity engaged in commercial, industrial, or professional activities. Businesses can be:
\begin{itemize}
    \item \textbf{For-profit entities} --- Organizations that operate to generate revenue and profit
    \item \textbf{Non-profit organizations} --- Organizations that operate to fulfill a charitable mission or further a social cause
\end{itemize}
\end{keyconceptbox}

Businesses range in scale from \textbf{sole proprietorships} to \textbf{international corporations} and can range in size from small to large.

\subsection{How Do You Start a Business?}

There are a number of steps you need to go through in order to start a business:

\begin{enumerate}[label=\textcolor{primaryblue}{\arabic*.}]
    \item \textbf{Conducting market research} --- Understanding your target market and competition
    \item \textbf{Developing a business plan} --- Creating a roadmap for your business
    \item \textbf{Seeking capital or other forms of funding} --- Securing financial resources
    \item \textbf{Choosing a location and business structure} --- Determining where and how to operate
    \item \textbf{Picking the right name} --- Selecting an appropriate business name
    \item \textbf{Submitting registration paperwork} --- Completing legal registration
    \item \textbf{Obtaining tax documents} --- Getting employer and taxpayer IDs
    \item \textbf{Pulling permits and licenses} --- Obtaining necessary legal permissions
    \item \textbf{Setting up a bank account} --- Facilitating everyday banking needs
\end{enumerate}

% ============================
% CHAPTER 2: NEEDS AND WANTS
% ============================
\section{Understanding Needs and Wants}

\begin{importantbox}
For the purpose of spending and saving money wisely, every person must know the difference between needs and wants.
\end{importantbox}

\subsection{What are Needs?}

\begin{definitionbox}
\textbf{Needs} are the basic requirements or necessities that are essential for your survival. In other words, they are the things you \textbf{have to have}, or \textbf{cannot do without}.
\end{definitionbox}

\subsubsection{Basic Needs for Human Survival}
\begin{itemize}[label=\textcolor{primaryblue}{\faIcon{check}}]
    \item \textbf{Oxygen} --- Essential for breathing
    \item \textbf{Water} --- Required for hydration
    \item \textbf{Food} --- Necessary for nutrition
    \item Basic health care products
    \item Clothing
    \item Tools to work
    \item Protective housing
\end{itemize}

\subsubsection{Types of Needs}

\begin{table}[h]
\centering
\begin{tabular}{|p{4cm}|p{9cm}|}
\hline
\rowcolor{primaryblue}
\textcolor{white}{\textbf{Type of Need}} & \textcolor{white}{\textbf{Description}} \\
\hline
\textbf{Physical Needs} & Basic or physiological needs such as food, water, and clothing \\
\hline
\textbf{Subjective Needs} & Psychological needs including self-esteem, sense of security, and love \\
\hline
\end{tabular}
\caption{Classification of Human Needs}
\end{table}

\subsection{What are Wants?}

\begin{definitionbox}
\textbf{Wants} are desires. Wants are the things that you would like to have but are \textbf{not essential} to your survival. Every person has their own set of desires or wants.
\end{definitionbox}

\begin{examplebox}
Examples of wants include:
\begin{itemize}
    \item A big house
    \item Fancy food
    \item Branded clothes
    \item A new car
    \item Trip to a foreign country
\end{itemize}
\end{examplebox}

\begin{keyconceptbox}
\textbf{Key Points about Wants:}
\begin{itemize}
    \item Wants can change over time as tastes and preferences change
    \item The inability to fulfill a want may not cause death or physical illness
    \item However, unfulfilled wants may result in mental distress
\end{itemize}
\end{keyconceptbox}

\subsection{Consumer Goods vs Producer Goods}

\begin{table}[h]
\centering
\begin{tabular}{|p{6cm}|p{6cm}|}
\hline
\rowcolor{primaryblue}
\textcolor{white}{\textbf{Consumer Goods}} & \textcolor{white}{\textbf{Producer Goods}} \\
\hline
Goods finally used by consumers to satisfy their wants & Goods used by producers to produce more goods \\
\hline
Directly satisfy human wants & Continue the process of production \\
\hline
\end{tabular}
\caption{Comparison of Consumer and Producer Goods}
\end{table}

% ============================
% CHAPTER 3: TYPES OF BUSINESS ORGANIZATIONS
% ============================
\section{Types of Business Organizations}

\subsection{Sole Trader}

\begin{definitionbox}
A \textbf{sole trader} is a single owner business, usually family orientated. A sole trader can easily control this type of business because it is small.
\end{definitionbox}

\begin{advantagebox}
\begin{itemize}
    \item Easy to control due to small size
    \item Limited capital required
    \item Few legal formalities
    \item Quick decision-making
\end{itemize}
\end{advantagebox}

\begin{disadvantagebox}
\begin{itemize}
    \item \textbf{Unlimited liability} --- If the business goes bankrupt, the owner has to forfeit personal possessions like land papers and personal car
\end{itemize}
\end{disadvantagebox}

\subsubsection{Sources of Fixed Capital for Sole Traders}
\begin{itemize}[label=\textcolor{primaryblue}{\faIcon{money-bill-wave}}]
    \item Personal savings of the sole trader
    \item Loans from banks
    \item Mortgages from building societies
    \item Ploughed-back profits
\end{itemize}

\subsection{Partnership}

\begin{definitionbox}
\textbf{Partnership} is a business between two to twenty persons. Each partner needs to sign a \textbf{Deed of Agreement} that gives details of:
\begin{itemize}
    \item Capital contributed by each partner
    \item Rate of interest to be paid on capital
    \item Salaries
\end{itemize}
\end{definitionbox}

\begin{advantagebox}
\begin{itemize}
    \item More capital available compared to a sole trader
    \item Partners can consult each other and make correct decisions
    \item Partnership can make use of individual skills
\end{itemize}
\end{advantagebox}

\begin{disadvantagebox}
\begin{itemize}
    \item Ordinary partners have \textbf{unlimited liability}
    \item Conflicts and disagreements between partners may occur
    \item Limited membership makes it difficult to raise large sums of capital
\end{itemize}
\end{disadvantagebox}

\subsubsection{Sources of Fixed Capital for Partnerships}
\begin{itemize}[label=\textcolor{primaryblue}{\faIcon{money-bill-wave}}]
    \item Personal savings of the partners
    \item Loans from banks
    \item Hire purchase and leasing
\end{itemize}

\subsection{Limited Companies}

\begin{keyconceptbox}
Limited companies have three key features:
\begin{enumerate}
    \item \textbf{Limited Liability} --- Shareholders only lose their investment if the company fails
    \item \textbf{Separate Legal Identity} --- The company is legally separate from its owners
    \item \textbf{Continuity} --- The company continues to exist even if ownership changes
\end{enumerate}
\end{keyconceptbox}

Companies must be filed with the \textbf{Registrar of Companies}. The shareholders elect a \textbf{Board of Directors}, usually at an annual general meeting, to run the company on their behalf.

\subsubsection{Sources of Long-Term Finance for Companies}
\begin{itemize}[label=\textcolor{primaryblue}{\faIcon{chart-line}}]
    \item Issuing debentures and shares
    \item Using ploughed-back profits
    \item Borrowing from banks
    \item Using hire purchase and leasing facilities
\end{itemize}

\subsection{Private Limited Companies (Ltd)}

\begin{definitionbox}
\textbf{Private Limited Companies} must use the words '\textbf{limited}' or '\textbf{Ltd}' after their name and cannot sell shares to the general public.
\end{definitionbox}

\begin{keyconceptbox}
\textbf{Key Requirements:}
\begin{itemize}
    \item One or more directors required
    \item Minimum of two shareholders
    \item Maximum number of shareholders is 50
    \item Cannot publicly trade shares
    \item Must be registered with the Registrar of Companies
    \item Accounts must be filed with the Registrar of Companies
\end{itemize}
\end{keyconceptbox}

\begin{examplebox}
Examples of Private Limited Companies in Bangladesh:
\begin{itemize}
    \item Elite Motors
    \item Unicode
    \item Progressive Packers and Shippers (Pvt) Limited
    \item Dell Service Bangladesh
\end{itemize}
\end{examplebox}

\subsection{Public Limited Companies (PLC)}

\begin{definitionbox}
\textbf{Public Limited Companies} must use the words '\textbf{public limited company}' or '\textbf{plc}' after their name.
\end{definitionbox}

\begin{keyconceptbox}
\textbf{Key Requirements:}
\begin{itemize}
    \item Minimum of two directors
    \item Minimum of two shareholders
    \item Members of the public can be invited to buy shares
    \item Shares are quoted on the stock exchange
\end{itemize}
\end{keyconceptbox}

\begin{advantagebox}
\begin{itemize}
    \item Capital is easy to raise due to:
    \begin{itemize}
        \item Limited liability for shareholders
        \item Public can buy shares
        \item Shares traded on stock exchange
    \end{itemize}
\end{itemize}
\end{advantagebox}

\begin{disadvantagebox}
\begin{itemize}
    \item Formation is costly in time and money
    \item Many documents must be submitted to the Registrar of Companies
    \item Coordination and communication problems might arise
\end{itemize}
\end{disadvantagebox}

\subsection{Incorporated vs Unincorporated Businesses}

\begin{table}[h]
\centering
\begin{tabular}{|p{6cm}|p{6cm}|}
\hline
\rowcolor{primaryblue}
\textcolor{white}{\textbf{Incorporated Businesses}} & \textcolor{white}{\textbf{Unincorporated Businesses}} \\
\hline
Limited liability & Unlimited liability \\
\hline
Separate legal personality & No separate legal identity \\
\hline
Can continue if bankruptcy occurs & May cease to exist \\
\hline
Must be registered under Companies Act & No registration required \\
\hline
Examples: ACI HealthCare Limited, United Hospital Limited & Examples: Sole traders, Partnerships \\
\hline
\end{tabular}
\caption{Comparison of Incorporated and Unincorporated Businesses}
\end{table}

\subsubsection{Important Documents for Incorporation}
\begin{itemize}[label=\textcolor{primaryblue}{\faIcon{file-alt}}]
    \item \textbf{Certificate of Incorporation} --- Issued by the Registrar of Companies after registration
    \item \textbf{Memorandum of Association} --- The document that sets up the company
    \item \textbf{Articles of Association} --- Sets out how the company is run, governed, and owned
\end{itemize}

\subsection{Cooperatives}

\begin{definitionbox}
\textbf{Cooperatives} are people-centred enterprises owned, controlled, and run by and for their members to realise their common economic, social, and cultural needs and aspirations.
\end{definitionbox}

\begin{keyconceptbox}
\textbf{Key Features of Cooperatives:}
\begin{itemize}
    \item Bring people together in a democratic and equal way
    \item Members vote on major decisions
    \item Board of directors elected from among members
\end{itemize}
\end{keyconceptbox}

\subsubsection{Types of Cooperatives}
\begin{itemize}[label=\textcolor{primaryblue}{\faIcon{users}}]
    \item \textbf{Consumers' Cooperative} --- Business owned by its customers (e.g., grocery stores)
    \item \textbf{Retailers' Cooperative} --- Members are retailers rather than consumers
\end{itemize}

% ============================
% CHAPTER 4: FRANCHISES
% ============================
\section{Franchise Business}

\begin{definitionbox}
A \textbf{franchise business} is a business in which the owners, or ``\textbf{franchisors}'', sell the rights to their business logo, name, and model to third party retail outlets, owned by independent, third party operators, called ``\textbf{franchisees}''.
\end{definitionbox}

\begin{keyconceptbox}
Franchises are an extremely common way of doing business around the world!
\end{keyconceptbox}

\subsection{The Role of the Franchisor}

\begin{itemize}[label=\textcolor{primaryblue}{\faIcon{building}}]
    \item Motivated to share their experience and know-how with franchisees
    \item Provide franchisees with tools needed to operate their businesses to brand standards
    \item Focused on ensuring each franchisee operates to system standards
\end{itemize}

\subsection{The Role of the Franchisee}

\begin{itemize}[label=\textcolor{primaryblue}{\faIcon{store}}]
    \item Ensure all employees are properly trained
    \item Ensure the franchise is properly staffed at all times
    \item Advertise and promote the franchise and its approved products
    \item Follow guidelines provided by the franchisor
\end{itemize}

% ============================
% CHAPTER 5: STAKEHOLDERS
% ============================
\section{Stakeholders}

\begin{definitionbox}
A \textbf{stakeholder} has a vested interest in a company and can either affect or be affected by a business's operations and performance.
\end{definitionbox}

\subsection{Types of Stakeholders}

\begin{table}[h]
\centering
\begin{tabular}{|p{6cm}|p{6cm}|}
\hline
\rowcolor{primaryblue}
\textcolor{white}{\textbf{Internal Stakeholders}} & \textcolor{white}{\textbf{External Stakeholders}} \\
\hline
People whose interest comes through a direct relationship & People who do not directly work with the company \\
\hline
Connected through employment, ownership, or investment & Affected by actions and outcomes of the business \\
\hline
Examples: Employees, Owners, Investors & Examples: Suppliers, Creditors, Public groups \\
\hline
\end{tabular}
\caption{Comparison of Internal and External Stakeholders}
\end{table}

\begin{keyconceptbox}
\textbf{Typical Stakeholders Include:}
\begin{itemize}[label=\textcolor{primaryblue}{\faIcon{user-tie}}]
    \item Investors
    \item Employees
    \item Customers
    \item Suppliers
    \item Communities
    \item Governments
    \item Trade associations
\end{itemize}
\end{keyconceptbox}

% ============================
% CHAPTER 6: INDUSTRY SECTORS
% ============================
\section{Primary, Secondary, and Tertiary Sectors}

\begin{keyconceptbox}
There are \textbf{three main types of industry} in which firms operate. These sectors form a \textbf{chain of production} which provides customers with finished goods or services.
\end{keyconceptbox}

\subsection{Primary Production (Extractive)}

\begin{definitionbox}
\textbf{Primary production} involves acquiring raw materials. This is sometimes known as \textbf{extractive production}.
\end{definitionbox}

\begin{examplebox}
Examples of Primary Production:
\begin{itemize}
    \item Metals and coal being \textbf{mined}
    \item Oil \textbf{drilled} from the ground
    \item Rubber \textbf{tapped} from trees
    \item Foodstuffs \textbf{farmed}
    \item Fish \textbf{trawled}
\end{itemize}
\end{examplebox}

\subsection{Secondary Production (Manufacturing)}

\begin{definitionbox}
\textbf{Secondary production} is the manufacturing and assembly process.
\end{definitionbox}

This involves:
\begin{itemize}[label=\textcolor{primaryblue}{\faIcon{industry}}]
    \item Converting raw materials into components (e.g., making plastics from oil)
    \item Assembling the product (e.g., building houses, bridges, and roads)
\end{itemize}

\subsection{Tertiary Production (Services)}

\begin{definitionbox}
\textbf{Tertiary production} refers to the commercial services that support the production and distribution process.
\end{definitionbox}

\begin{examplebox}
Examples of Tertiary Services:
\begin{itemize}[columns=2]
    \item Insurance
    \item Transport
    \item Advertising
    \item Warehousing
    \item Teaching
    \item Health care
\end{itemize}
\end{examplebox}

\begin{figure}[h]
\centering
\begin{tikzpicture}[node distance=2.5cm]
    \node[draw, rounded corners, fill=primaryblue, text=white, minimum width=3cm, minimum height=1cm] (primary) {\textbf{Primary Sector}};
    \node[draw, rounded corners, fill=secondaryblue, text=white, minimum width=3cm, minimum height=1cm, right of=primary, xshift=2cm] (secondary) {\textbf{Secondary Sector}};
    \node[draw, rounded corners, fill=accentgreen, text=white, minimum width=3cm, minimum height=1cm, right of=secondary, xshift=2cm] (tertiary) {\textbf{Tertiary Sector}};
    
    \draw[->, line width=2pt, darkgray] (primary) -- (secondary);
    \draw[->, line width=2pt, darkgray] (secondary) -- (tertiary);
\end{tikzpicture}
\caption{Chain of Production}
\end{figure}

% ============================
% CHAPTER 7: TRAINING
% ============================
\section{Training in Business}

\begin{keyconceptbox}
\textbf{Training} is basically a teaching and learning activity carried on to enhance the skills, knowledge, and abilities of the employees working in the organization.
\end{keyconceptbox}

\subsection{Induction Training}

\begin{definitionbox}
\textbf{Induction training} enables a new recruit to become productive as quickly as possible. It can avoid costly mistakes by recruits not knowing the procedures or techniques of their new jobs.
\end{definitionbox}

The length of induction training will vary based on:
\begin{itemize}[label=\textcolor{primaryblue}{\faIcon{clock}}]
    \item Complexity of the job
    \item Size of the business
    \item Level or position of the job within the business
\end{itemize}

\subsubsection{Areas Included in Induction Training}
\begin{itemize}[label=\textcolor{primaryblue}{\faIcon{graduation-cap}}]
    \item Learning about the duties of the job
    \item Meeting new colleagues
    \item Seeing the layout of the premises
    \item Learning the values and aims of the business
    \item Learning about the internal workings and policies of the business
\end{itemize}

\subsection{On-the-Job Training}

\begin{definitionbox}
With \textbf{on-the-job training}, employees receive training whilst remaining in the workplace.
\end{definitionbox}

\subsubsection{Main Methods of On-the-Job Training}

\begin{enumerate}[label=\textcolor{primaryblue}{\arabic*.}]
    \item \textbf{Demonstration/Instruction} --- Showing the trainee how to do the job
    \item \textbf{Coaching} --- A more intensive method involving a close working relationship between an experienced employee and the trainee
    \item \textbf{Job Rotation} --- The trainee is given several jobs in succession to gain experience of a wide range of activities (e.g., a graduate management trainee might spend periods in several different departments)
    \item \textbf{Projects} --- Employees join a project team, giving them exposure to other parts of the business and allowing them to take part in new activities. Most successful project teams are ``multi-disciplinary''
\end{enumerate}

\begin{table}[h]
\centering
\begin{tabular}{|p{6cm}|p{6cm}|}
\hline
\rowcolor{primaryblue}
\multicolumn{2}{|c|}{\textcolor{white}{\textbf{On-the-Job Training}}} \\
\hline
\rowcolor{accentgreen!20}
\textbf{Advantages} & \rowcolor{accentred!20}\textbf{Disadvantages} \\
\hline
Generally most cost-effective & Quality depends on ability of trainer and time available \\
\hline
Employees are actually productive & Bad habits might be passed on \\
\hline
Opportunity to learn whilst doing & Learning environment may not be conducive \\
\hline
Training alongside real colleagues & Potential disruption to production \\
\hline
\end{tabular}
\caption{Advantages and Disadvantages of On-the-Job Training}
\end{table}

\subsection{Off-the-Job Training}

\begin{definitionbox}
\textbf{Off-the-job training} occurs when employees are taken away from their place of work to be trained.
\end{definitionbox}

\subsubsection{Common Methods of Off-the-Job Training}
\begin{itemize}[label=\textcolor{primaryblue}{\faIcon{university}}]
    \item \textbf{Day release} --- Employee takes time off work to attend a local college or training centre
    \item \textbf{Distance learning/evening classes}
    \item \textbf{Block release courses} --- May involve several weeks at a local college
    \item \textbf{Sandwich courses} --- Employee spends a longer period at college (e.g., six months) before returning to work
    \item \textbf{Sponsored courses in higher education}
    \item \textbf{Self-study, computer-based training}
\end{itemize}

\begin{table}[h]
\centering
\begin{tabular}{|p{6cm}|p{6cm}|}
\hline
\rowcolor{primaryblue}
\multicolumn{2}{|c|}{\textcolor{white}{\textbf{Off-the-Job Training}}} \\
\hline
\rowcolor{accentgreen!20}
\textbf{Advantages} & \rowcolor{accentred!20}\textbf{Disadvantages} \\
\hline
A wider range of skills or qualifications can be obtained & More expensive (transport and accommodation) \\
\hline
Can learn from outside specialists or experts & Lost working time and potential output \\
\hline
Employees can be more confident when starting job & New employees may still need induction training \\
\hline
 & Employees may leave for better jobs with new skills \\
\hline
\end{tabular}
\caption{Advantages and Disadvantages of Off-the-Job Training}
\end{table}

% ============================
% CHAPTER 8: RECRUITMENT AND SELECTION
% ============================
\section{Recruitment and Selection}

\begin{definitionbox}
\textbf{Recruitment} is the process from identifying that the business needs to employ someone up to the point where applications have arrived at the business.
\end{definitionbox}

\subsection{When Does a Vacancy Arise?}

A vacancy arises when an employee:
\begin{itemize}[label=\textcolor{primaryblue}{\faIcon{user-minus}}]
    \item Resigns from a job
    \item Is dismissed by the management
\end{itemize}

\subsection{The Recruitment Process}

\subsubsection{Step 1: Job Analysis}

\begin{keyconceptbox}
A \textbf{job analysis} identifies and records the tasks and responsibilities relating to the job. It will tell the managers what the job post is for.
\end{keyconceptbox}

\subsubsection{Step 2: Job Description}

\begin{definitionbox}
A \textbf{job description} outlines the responsibilities and duties to be carried out by someone employed to do the job.
\end{definitionbox}

A job description includes:
\begin{itemize}[label=\textcolor{primaryblue}{\faIcon{file-contract}}]
    \item Conditions of employment (salary, working hours, pension scheme)
    \item Training offered
    \item Opportunities for promotion
\end{itemize}

This is given to all prospective candidates so they know exactly what they will be required and expected to do.

\subsubsection{Step 3: Job Specification}

\begin{definitionbox}
A \textbf{job specification} is a document that outlines the requirements, qualifications, expertise, skills, physical/personal characteristics etc. required by an employee to be able to take up the job.
\end{definitionbox}

\subsection{Advertising the Vacancy}

\subsubsection{Internal Recruitment}

\begin{definitionbox}
\textbf{Internal recruitment} is when a vacancy is filled by an \textbf{existing employee} of the business.
\end{definitionbox}

\begin{advantagebox}
\begin{itemize}
    \item Saves time and money --- no need for advertising and interviewing
    \item Person already known to business
    \item Person knows business's ways of working
    \item Motivating for other employees to see colleagues being promoted
\end{itemize}
\end{advantagebox}

\begin{disadvantagebox}
\begin{itemize}
    \item No new skills and experience coming into the business
    \item Jealousy among workers
\end{itemize}
\end{disadvantagebox}

\subsubsection{External Recruitment}

\begin{definitionbox}
\textbf{External recruitment} is when a vacancy is filled by someone who is \textbf{not an existing employee} and will be new to the business.
\end{definitionbox}

\textbf{Methods of External Advertising:}
\begin{itemize}[label=\textcolor{primaryblue}{\faIcon{newspaper}}]
    \item Local/national newspapers
    \item Specialist magazines and journals
    \item Job centres run by the government
    \item Recruitment agencies
\end{itemize}

\begin{importantbox}
When advertising a job, the business needs to consider:
\begin{itemize}
    \item What should be included in the advertisement
    \item Where it should be advertised
    \item How much it will cost
    \item Whether it will be cost-effective
\end{itemize}
\end{importantbox}

\subsubsection{Curriculum Vitae (CV)}

When a person is interested in a job, they should apply by sending in a \textbf{curriculum vitae (CV)} or resume, which includes:
\begin{itemize}[label=\textcolor{primaryblue}{\faIcon{user-circle}}]
    \item Qualifications
    \item Experience
    \item Qualities and skills
    \item Statements about why they want the job
    \item Why they feel suitable for the job
\end{itemize}

\subsection{Selection Process}

\begin{keyconceptbox}
Applicants who are \textbf{shortlisted} will be interviewed by the H.R. manager. They will also call up the \textbf{referee} provided by the applicant (a referee could be the previous employer or colleagues who can give a confidential opinion about the applicant's reliability, honesty, and suitability for the job).
\end{keyconceptbox}

\subsubsection{What Interviews Assess}
\begin{itemize}[label=\textcolor{primaryblue}{\faIcon{clipboard-check}}]
    \item The applicant's ability to do the job
    \item Personal qualities of the applicant
    \item Character and personality of applicant
\end{itemize}

\subsubsection{Selection Tests}
\begin{itemize}[label=\textcolor{primaryblue}{\faIcon{tasks}}]
    \item \textbf{Skills tests} --- Ability to do the job
    \item \textbf{Aptitude tests} --- Candidate's potential to gain additional skills
    \item \textbf{Personality tests} --- What kind of personality the candidate has
    \item \textbf{Group situation tests} --- How they manage and work in teams
\end{itemize}

\subsection{Contract of Employment}

\begin{definitionbox}
The \textbf{contract of employment} is a legal agreement between the employer and the employee listing the rights and responsibilities of workers.
\end{definitionbox}

\textbf{A Contract of Employment Includes:}
\begin{itemize}[label=\textcolor{primaryblue}{\faIcon{file-signature}}]
    \item Name of employer and employee
    \item Job title
    \item Date when employment will begin
    \item Hours to work
    \item Rate of pay and other benefits
    \item When payment is made
    \item Holiday entitlement
    \item Notice period required to terminate employment
\end{itemize}

\subsection{Part-Time vs Full-Time Employment}

\begin{keyconceptbox}
\begin{itemize}
    \item \textbf{Part-time employment:} 1 to 30--35 hours a week
    \item \textbf{Full-time employment:} 35 hours or more a week
\end{itemize}
\end{keyconceptbox}

\begin{table}[h]
\centering
\begin{tabular}{|p{6cm}|p{6cm}|}
\hline
\rowcolor{primaryblue}
\multicolumn{2}{|c|}{\textcolor{white}{\textbf{Part-Time Employment}}} \\
\hline
\rowcolor{accentgreen!20}
\textbf{Advantages to Employer} & \rowcolor{accentred!20}\textbf{Disadvantages to Employer} \\
\hline
More flexible hours of work & Less likely to be trained (workers see job as temporary) \\
\hline
Easier to ask employees to work at busy times & Takes longer to recruit two part-time workers than one full-time worker \\
\hline
Easier to extend business opening hours & Can be less committed to the business \\
\hline
Workers accept lower pay for lesser hours & Less likely to be promoted \\
\hline
Less expensive than full-time workers & More difficult to communicate when not in work \\
\hline
\end{tabular}
\caption{Part-Time Employment: Advantages and Disadvantages}
\end{table}

% ============================
% CHAPTER 9: PRIVATISATION
% ============================
\section{Privatisation}

\begin{definitionbox}
\textbf{Privatisation} is the process of transferring an enterprise or industry from the \textbf{public sector} to the \textbf{private sector}.
\end{definitionbox}

\begin{examplebox}
Examples of privatised services:
\begin{itemize}
    \item Airport operation
    \item Data processing
    \item Vehicle maintenance
\end{itemize}
\end{examplebox}

\subsection{Potential Benefits of Privatisation}

\begin{advantagebox}
\begin{enumerate}[label=\textcolor{accentgreen}{\arabic*.}]
    \item \textbf{Improved Efficiency} --- Private companies have a profit incentive to cut costs and be more efficient
    
    \item \textbf{Lack of Political Interference} --- Governments are motivated by political pressures rather than sound economic and business sense
    
    \item \textbf{Shareholder Pressure} --- A private firm has pressure from shareholders to perform efficiently. If inefficient, the firm could be subject to a takeover
    
    \item \textbf{Increased Competition} --- Privatisation often occurs alongside deregulation, allowing more firms to enter the industry and increase market competitiveness
\end{enumerate}
\end{advantagebox}

\subsection{Disadvantages of Privatisation}

\begin{disadvantagebox}
\begin{enumerate}[label=\textcolor{accentred}{\arabic*.}]
    \item \textbf{Natural Monopoly} --- When the most efficient number of firms in an industry is one (e.g., tap water has significant fixed costs, leaving no scope for competition)
    
    \item \textbf{Public Interest} --- Industries like health care, education, and public transport perform important public services where profit shouldn't be the primary objective
    
    \item \textbf{Government Loses Dividends} --- Many privatised companies are profitable, meaning the government misses out on dividends that go to wealthy shareholders
    
    \item \textbf{Regulation of Private Monopolies} --- Privatisation creates private monopolies (like water and rail companies) that need regulating to prevent abuse of monopoly power
\end{enumerate}
\end{disadvantagebox}

% ============================
% CHAPTER 10: RETAILERS
% ============================
\section{Retailers}

\begin{definitionbox}
\textbf{Retailers} are the final link in the channel of distribution. A retailer purchases in bulk from the wholesalers and sells the products to the customers in small quantities.
\end{definitionbox}

\begin{keyconceptbox}
The aim of a retailer is to achieve \textbf{maximum customer satisfaction} by exceeding their expectations and delivering exceptional services.
\end{keyconceptbox}

\subsection{Services Provided by Retailers}

\begin{itemize}[label=\textcolor{primaryblue}{\faIcon{shopping-cart}}]
    \item To provide the goods that consumers want
    \item To provide goods in suitable quantity for consumers
    \item To provide the goods when customers want them
    \item Providing credit facilities
    \item Advertising and window displays
    \item Introducing new products
\end{itemize}

% ============================
% CHAPTER 11: INTERNATIONAL TRADE
% ============================
\section{International Trade}

\begin{definitionbox}
\textbf{International trade} is the exchange of services, goods, and capital among various countries and regions, without much hindrance. International trade accounts for a good part of a country's gross domestic product.
\end{definitionbox}

\begin{keyconceptbox}
International trade has flourished over the years due to:
\begin{itemize}
    \item Modern production techniques
    \item Highly advanced transportation systems
    \item Transnational corporations
    \item Outsourcing of manufacturing and services
    \item Rapid industrialisation
\end{itemize}
\end{keyconceptbox}

\subsection{Benefits of International Trade}

\begin{advantagebox}
\begin{itemize}[label=\textcolor{accentgreen}{\faIcon{globe}}]
    \item Enhancing the domestic competitiveness
    \item Taking advantage of international trade technology
    \item Increasing sales and profits
    \item Extending sales potential of existing products
    \item Maintaining cost competitiveness in the domestic market
    \item Enhancing potential for expansion of business
    \item Gaining a global market share
    \item Reducing dependence on existing markets
    \item Stabilising seasonal market fluctuations
\end{itemize}
\end{advantagebox}

\subsection{Imports and Exports}

\begin{table}[h]
\centering
\begin{tabular}{|p{6cm}|p{6cm}|}
\hline
\rowcolor{primaryblue}
\textcolor{white}{\textbf{Imports}} & \textcolor{white}{\textbf{Exports}} \\
\hline
Foreign goods entering a domestic economy & Home products moving to foreign countries \\
\hline
\multicolumn{2}{|c|}{Both are crucial for business effectiveness and competition} \\
\hline
\end{tabular}
\caption{Imports vs Exports}
\end{table}

\subsection{Visible and Invisible Trade}

\begin{keyconceptbox}
\begin{itemize}
    \item \textbf{Visible trade} --- Trade in \textbf{goods} between countries
    \item \textbf{Invisible trade} --- Trade in \textbf{services} between countries
\end{itemize}
\end{keyconceptbox}

\subsubsection{Balance of Trade Calculations}

\begin{tcolorbox}[colback=lightblue, colframe=primaryblue, title=\textbf{Balance of Trade Formulas}]
\textbf{Balance of Trade (Visible Balance):}
\[ \text{Visible Exports} - \text{Visible Imports} \]

\textbf{Invisible Balance:}
\[ \text{Invisible Exports} - \text{Invisible Imports} \]

\textbf{Balance of Payments on Current Account:}
\[ (\text{Visible Exports} + \text{Invisible Exports}) - (\text{Visible Imports} + \text{Invisible Imports}) \]

\begin{itemize}
    \item \textbf{Deficit} = Loss in the balance of payments
    \item \textbf{Surplus} = Profit in the balance of payments
\end{itemize}
\end{tcolorbox}

\subsection{Trade Blocs}

\begin{definitionbox}
A \textbf{trade bloc} is a type of intergovernmental agreement, often part of a regional intergovernmental organisation, where barriers to trade (tariffs and others) are reduced or eliminated among the participating states.
\end{definitionbox}

\subsubsection{Types of Trade Blocs}
\begin{enumerate}[label=\textcolor{primaryblue}{\arabic*.}]
    \item Preferential trading areas
    \item Free-trade areas
    \item Customs unions
    \item Common markets
    \item Economic and monetary unions
\end{enumerate}

\begin{examplebox}
Examples of Trade Blocs:
\begin{itemize}
    \item North American Free Trade Agreement (NAFTA)
    \item European Union (EU)
\end{itemize}
\end{examplebox}

\subsubsection{Benefits of Trade Blocs}
\begin{itemize}[label=\textcolor{primaryblue}{\faIcon{handshake}}]
    \item Removal of tariffs within trading bloc countries
    \item Common external tariff
    \item Free movement of labour, capital, and goods
    \item Better scope for improvement and market share
\end{itemize}

\subsection{Difficulties Facing Exporters}

\begin{importantbox}
Common problems faced by exporters worldwide:
\begin{itemize}[label=\textcolor{accentorange}{\faIcon{exclamation-circle}}]
    \item Language barriers
    \item Distance
    \item Methods of payment
    \item Differing consumer wants
    \item Documentation requirements
    \item Currency exchange issues
\end{itemize}
\end{importantbox}

\subsection{Multinational Corporations}

\begin{definitionbox}
\textbf{Multinationals} are companies with significant production operations in at least two countries.
\end{definitionbox}

\begin{examplebox}
Examples: Shell, DHL
\end{examplebox}

\begin{advantagebox}
\textbf{Benefits of Multinationals:}
\begin{itemize}
    \item Employment increases in the country
    \item Competition benefits as businesses can expand
    \item Training facilities increase for local labour force
    \item World-class goods become cheaper
    \item Many countries can enjoy products from different countries
\end{itemize}
\end{advantagebox}

\begin{disadvantagebox}
\textbf{Drawbacks of Multinationals:}
\begin{itemize}
    \item \textbf{Outflows of money} --- Home country takes back most of their finance
    \item \textbf{Worsening balance of payments} --- Multinational remits (sends) profits to its home country
\end{itemize}
\end{disadvantagebox}

% ============================
% SUMMARY
% ============================
\section{Summary}

\begin{tcolorbox}[colback=lightgray, colframe=darkgray, title=\textbf{\faIcon{list-check} Key Takeaways}]
\begin{enumerate}[label=\textcolor{primaryblue}{\arabic*.}]
    \item \textbf{Business Studies} combines accountancy, finance, marketing, HR, and operations
    \item Understanding \textbf{needs vs wants} is essential for wise financial decisions
    \item \textbf{Business types} include sole traders, partnerships, and limited companies
    \item \textbf{Incorporated businesses} have limited liability and separate legal identity
    \item \textbf{Stakeholders} can be internal or external to an organization
    \item The \textbf{three sectors} (primary, secondary, tertiary) form the chain of production
    \item \textbf{Training} enhances employee skills, knowledge, and abilities
    \item \textbf{Recruitment} involves job analysis, description, specification, and selection
    \item \textbf{Privatisation} transfers enterprises from public to private sector
    \item \textbf{International trade} facilitates global exchange of goods, services, and capital
\end{enumerate}
\end{tcolorbox}

\end{document}
